\documentclass[10pt,a4paper]{article}

% =========================
% Page + Font
% =========================
\usepackage[a4paper,left=1.6cm,right=1.6cm,top=1.3cm,bottom=1.3cm]{geometry}
\usepackage[T1]{fontenc}
\usepackage[utf8]{inputenc}
\usepackage{helvet}
\renewcommand{\familydefault}{\sfdefault}
\usepackage{microtype}
\usepackage{setspace}
\usepackage{enumitem}

% =========================
% Tables + Colors + Figures
% =========================
\usepackage[table]{xcolor}
\usepackage{colortbl}
\usepackage{array}
\usepackage{tabularx}
\usepackage{booktabs}
\usepackage{makecell}
\usepackage{graphicx}
\usepackage{float}
\usepackage{caption}
\usepackage{ragged2e}
\usepackage{verbatim}

% =========================
% General formatting
% =========================
\pagestyle{empty}
\setlength{\parindent}{0pt}
\setlength{\parskip}{2pt}
\setlength{\tabcolsep}{4pt}
\renewcommand{\arraystretch}{1.18}
\captionsetup{font=footnotesize}
\setlist[itemize]{leftmargin=1.2em,itemsep=1pt,topsep=2pt}
\setlist[enumerate]{leftmargin=1.2em,itemsep=1pt,topsep=2pt}

% =========================
% Palette (matched to template)
% =========================
\definecolor{ReportBlue}{HTML}{0D2C54}
\definecolor{LightBlue}{HTML}{1E5799}
\definecolor{LightGray}{HTML}{F4F6F9}
\definecolor{MidGray}{HTML}{C8CDD6}
\definecolor{GoodGreen}{HTML}{1E7D46}
\definecolor{MetricRed}{HTML}{C0392B}
\definecolor{CaptionGray}{HTML}{666666}
\definecolor{BodyGray}{HTML}{333333}
\definecolor{FooterGray}{HTML}{999999}
\definecolor{KPIBlueBG}{HTML}{EAF0FB}
\definecolor{KPIGreenBG}{HTML}{E8F5EE}

% =========================
% Reusable macros
% =========================
\newcommand{\SectionBar}[1]{%
  \vspace{0.35em}
  \noindent\colorbox{ReportBlue}{%
    \parbox{\dimexpr\linewidth-2\fboxsep\relax}{%
      \vspace{2.5pt}\textcolor{white}{\bfseries\footnotesize #1}\vspace{2.5pt}
    }%
  }%
  \vspace{0.35em}
}

% =========================
% Report variables
% =========================
\newcommand{\ReportTitle}{Comprehensive Project Report}
\newcommand{\ReportSubtitle}{Multi-Dimensional Return Forecasting and Portfolio Management for Indian Stocks | DA6701 Assignment 2 | Train: Jan 1, 2020 -- Sep 30, 2025 | Test: Oct 1, 2025 -- Dec 31, 2025}

\newcommand{\KPIOne}{0.2113}
\newcommand{\KPITwo}{7.33\%}
\newcommand{\KPIThree}{-2.52\%}
\newcommand{\KPIFour}{51.67\%}
\newcommand{\KPIFive}{26.39\%}
\newcommand{\KPISix}{0.000049}

\newcommand{\WeightsFig}{weights_predicted_return.png}
\newcommand{\EquityCurveFig}{equity_curve_predicted_return.png}

\newcommand{\FooterText}{Tools: Python 3, yfinance, requests, beautifulsoup4, newsapi-python, PyTorch, scipy.optimize (SLSQP), pandas, numpy, scikit-learn, transformers (FinBERT), matplotlib, seaborn, joblib, tqdm, jupyter}

\begin{document}
\color{BodyGray}

% =========================
% Title
% =========================
\begin{center}
  {\fontsize{15}{18}\selectfont\textbf{\textcolor{ReportBlue}{\ReportTitle}}}\par
  \vspace{1pt}
  {\fontsize{8}{10}\selectfont\textcolor{CaptionGray}{\ReportSubtitle}}
\end{center}
\vspace{4pt}
\noindent\textcolor{ReportBlue}{\rule{\linewidth}{2pt}}
\vspace{5pt}

% =========================
% KPI Banner
% =========================
\noindent
\begin{minipage}[t]{0.496\linewidth}
  \arrayrulecolor{LightBlue}
  \setlength{\tabcolsep}{3pt}
  \begin{tabular}{|>{\centering\arraybackslash}p{0.30\linewidth}|>{\centering\arraybackslash}p{0.30\linewidth}|>{\centering\arraybackslash}p{0.30\linewidth}|}
    \hline
    \cellcolor{KPIBlueBG}{\fontsize{13}{15}\selectfont\textbf{\textcolor{GoodGreen}{\KPIOne}}} &
    \cellcolor{KPIBlueBG}{\fontsize{13}{15}\selectfont\textbf{\textcolor{LightBlue}{\KPITwo}}} &
    \cellcolor{KPIBlueBG}{\fontsize{13}{15}\selectfont\textbf{\textcolor{MetricRed}{\KPIThree}}} \\
    \cellcolor{KPIBlueBG}{\fontsize{7}{9}\selectfont\textcolor{CaptionGray}{Sharpe Ratio}} &
    \cellcolor{KPIBlueBG}{\fontsize{7}{9}\selectfont\textcolor{CaptionGray}{Total Return}} &
    \cellcolor{KPIBlueBG}{\fontsize{7}{9}\selectfont\textcolor{CaptionGray}{Max Drawdown}} \\
    \hline
  \end{tabular}
\end{minipage}
\hfill
\begin{minipage}[t]{0.496\linewidth}
  \arrayrulecolor{GoodGreen}
  \setlength{\tabcolsep}{3pt}
  \begin{tabular}{|>{\centering\arraybackslash}p{0.30\linewidth}|>{\centering\arraybackslash}p{0.30\linewidth}|>{\centering\arraybackslash}p{0.30\linewidth}|}
    \hline
    \cellcolor{KPIGreenBG}{\fontsize{13}{15}\selectfont\textbf{\textcolor{LightBlue}{\KPIFour}}} &
    \cellcolor{KPIGreenBG}{\fontsize{13}{15}\selectfont\textbf{\textcolor{MetricRed}{\KPIFive}}} &
    \cellcolor{KPIGreenBG}{\fontsize{13}{15}\selectfont\textbf{\textcolor{GoodGreen}{\KPISix}}} \\
    \cellcolor{KPIGreenBG}{\fontsize{7}{9}\selectfont\textcolor{CaptionGray}{Portfolio Hit Ratio}} &
    \cellcolor{KPIGreenBG}{\fontsize{7}{9}\selectfont\textcolor{CaptionGray}{Directional Accuracy}} &
    \cellcolor{KPIGreenBG}{\fontsize{7}{9}\selectfont\textcolor{CaptionGray}{Best Model MSE (HDFCBANK)}} \\
    \hline
  \end{tabular}
\end{minipage}

\vspace{6pt}

% =========================
% Section 1
% =========================
\SectionBar{1 \;|\; Data and Feature Engineering}
{\fontsize{8}{10.3}\selectfont\justifying
\textbf{Data sources and collection.} The project combines four aligned daily data streams for 6 Indian stocks (RELIANCE, HDFCBANK, INFY, M\&M, BHARTIARTL, HUL) over Jan 1, 2020 to Dec 31, 2025.

\begin{tabularx}{\linewidth}{|>{\RaggedRight\arraybackslash}p{2.9cm}|>{\RaggedRight\arraybackslash}p{4.5cm}|X|}
\hline
\rowcolor{ReportBlue}
\textcolor{white}{\textbf{Data Type}} & \textcolor{white}{\textbf{Coverage}} & \textcolor{white}{\textbf{Source and Notes}} \\
\hline
Market Data (OHLCV) & Daily adjusted prices and volume (6 tickers) & Yahoo Finance via \texttt{yfinance}; period: Jan 1, 2020 -- Dec 31, 2025 \\
\hline
Fundamental Data & Quarterly P/E, Debt-to-Equity, ROE, EPS; mapped to daily & Scraped from screener.in and aligned by forward-fill \\
\hline
Macro Indicators & USD-INR, Brent oil, NIFTY 50, US 10Y (proxy), VIX & Yahoo Finance via \texttt{yfinance}; lagged to avoid look-ahead leakage \\
\hline
Sentiment Data & Daily polarity from financial headlines & NewsAPI headline retrieval + FinBERT transformer scoring \\
\hline
\end{tabularx}

\textbf{Feature engineering pipeline.}
\begin{enumerate}
  \item Returns are computed as log returns: $\ln(\text{Close}_t/\text{Close}_{t-1})$.
  \item Lagged features include: returns (t-1 to t-5), volume (t-1) with 10-day MA, fundamentals (t-1), macro (t-1), sentiment (t-1).
  \item Technical indicators include SMA(10), SMA(20), RSI(14), MACD line, signal line, and histogram.
  \item RobustScaler is fit on training data only and applied to both train/test to reduce outlier sensitivity.
\end{enumerate}

\textbf{Data split.} Training period: Jan 1, 2020 -- Sep 30, 2025 ($\sim$5.75 years). Testing period: Oct 1, 2025 -- Dec 31, 2025 (3 months forward test).
}

% =========================
% Section 2
% =========================
\SectionBar{2 \;|\; Model Architecture and Validation}
{\fontsize{8}{10.3}\selectfont\justifying
\textbf{Model architecture (per ticker).}
\begin{itemize}
  \item Type: Long Short-Term Memory (LSTM) recurrent neural network.
  \item Input: 30-day feature sequences.
  \item Core layers: 2 LSTM layers, 64 hidden units each, with dropout 0.3.
  \item Dense head: 32-unit fully connected layer with ReLU.
  \item Output: one unit predicting next-day log return.
\end{itemize}

\textbf{Hyperparameters and optimization.}
\begin{itemize}
  \item Optimizer: Adam with learning rate 0.001.
  \item Scheduler: ReduceLROnPlateau.
  \item Batch size: 32.
  \item Epoch budget: up to 30 epochs.
  \item Early stopping: patience = 5.
\end{itemize}

\textbf{Regularization and training controls.}
\begin{enumerate}
  \item Dropout (0.3) in recurrent and dense blocks.
  \item L2 weight decay ($1\times10^{-5}$).
  \item Gradient clipping with \texttt{max\_norm=1.0}.
  \item Early stopping on validation loss.
\end{enumerate}

\textbf{Validation strategy.} 5-fold walk-forward time-series validation with expanding windows was used to preserve chronology and avoid leakage. K-fold shuffle validation was not used.

\textbf{Training process.} Six independent models were trained (one per ticker) and selected using validation MSE across walk-forward folds.
}

% =========================
% Section 3
% =========================
\SectionBar{3 \;|\; Feature Importance and Interpretation}
{\fontsize{8}{10.3}\selectfont\justifying
Ablation and model diagnostics indicate the following importance order:
\begin{enumerate}
  \item Lagged returns (t-1 to t-5): strongest predictors for momentum/mean-reversion structure.
  \item Technical indicators (RSI, MACD): encode trend and short-horizon market behavior.
  \item Macro indicators (USD-INR, oil): capture exogenous macro pressure.
  \item Volume features: provide participation and liquidity signals.
  \item Sentiment scores: represent news-driven effects.
  \item Fundamental metrics (P/E, ROE): add slower-moving valuation context.
\end{enumerate}

Importance varies across symbols: RELIANCE and HDFCBANK are more macro-sensitive, while INFY shows stronger dependence on technical patterns.
}

% =========================
% Section 4
% =========================
\SectionBar{4 \;|\; Individual Model Results (Test: Oct--Dec 2025)}

{\fontsize{7.7}{9.6}\selectfont
\arrayrulecolor{MidGray}
\rowcolors{2}{LightGray}{white}
\begin{tabular}{|>{\centering\arraybackslash}p{2.15cm}|>{\centering\arraybackslash}p{2.15cm}|>{\centering\arraybackslash}p{2.15cm}|>{\centering\arraybackslash}p{2.15cm}|>{\centering\arraybackslash}p{2.25cm}|}
  \hline
  \rowcolor{ReportBlue}
  \textcolor{white}{\textbf{Ticker}} &
  \textcolor{white}{\textbf{MSE}} &
  \textcolor{white}{\textbf{MAE}} &
  \textcolor{white}{\textbf{Hit Ratio}} &
  \textcolor{white}{\textbf{Directional Acc}} \\
  \hline
  RELIANCE   & 0.000089 & 0.006737 & 0.317 & 0.264 \\
  HDFCBANK   & 0.000049 & 0.005473 & 0.283 & 0.264 \\
  INFY       & 0.000168 & 0.010058 & 0.267 & 0.264 \\
  M\&M       & 0.000112 & 0.008278 & 0.283 & 0.264 \\
  BHARTIARTL & 0.000150 & 0.009548 & 0.217 & 0.264 \\
  HUL        & 0.000115 & 0.007170 & 0.250 & 0.264 \\
  \hline
\end{tabular}
}

{\fontsize{8}{10.3}\selectfont\justifying
\textbf{Key observations.}
\begin{itemize}
  \item Cross-symbol robustness: all models achieved test MSE below 0.0002.
  \item Best test MSE: HDFCBANK (0.000049), with RELIANCE close behind.
  \item Hit ratios range from 0.217 to 0.317, with RELIANCE strongest in this set.
  \item Early stopping (patience = 5) was triggered to limit overfitting.
  \item Walk-forward validation MSE range: 0.000189 to 0.000463.
\end{itemize}
}

% =========================
% Section 5
% =========================
\SectionBar{5 \;|\; Portfolio Performance Metrics and Visual Diagnostics}

{\fontsize{7.4}{9.2}\selectfont
\arrayrulecolor{MidGray}
\rowcolors{2}{LightGray}{white}
\begin{tabular}{|>{\centering\arraybackslash}p{2.2cm}|>{\centering\arraybackslash}p{1.8cm}|>{\RaggedRight\arraybackslash}p{5.1cm}|}
  \hline
  \rowcolor{ReportBlue}
  \textcolor{white}{\textbf{Ticker}} &
  \textcolor{white}{\textbf{Weight}} &
  \textcolor{white}{\textbf{Rationale}} \\
  \hline
  RELIANCE   & 19.75\% & Moderate allocation, diversified \\
  HDFCBANK   & 20.38\% & Stable financial profile, strong balance \\
  INFY       & 28.43\% & Tech exposure, largest allocation \\
  M\&M       & 3.13\%  & Higher volatility, minimized allocation \\
  BHARTIARTL & 14.74\% & Telecom stability \\
  HUL        & 13.58\% & FMCG defensive characteristics \\
  \hline
\end{tabular}
}

{\fontsize{7.4}{9.2}\selectfont
\arrayrulecolor{MidGray}
\rowcolors{2}{LightGray}{white}
\begin{tabular}{|>{\RaggedRight\arraybackslash}p{3.8cm}|>{\centering\arraybackslash}p{2.1cm}|>{\RaggedRight\arraybackslash}p{3.6cm}|}
  \hline
  \rowcolor{ReportBlue}
  \textcolor{white}{\textbf{Metric}} &
  \textcolor{white}{\textbf{Value}} &
  \textcolor{white}{\textbf{Interpretation}} \\
  \hline
  Sharpe Ratio & 0.2113 & Risk-adjusted return (rf = 0 assumed) \\
  Maximum Drawdown & -2.52\% & Largest peak-to-trough decline \\
  Hit Ratio & 51.67\% & Share of days with positive returns \\
  Directional Accuracy & 26.39\% & Correct sign prediction rate \\
  Total Return & 7.33\% & Cumulative return over 3 months \\
  Mean Daily Return & 0.1196\% & Average daily return \\
  Daily Volatility & 0.5661\% & Standard deviation of returns \\
  \hline
\end{tabular}
}

\vspace{2pt}
\noindent
\begin{minipage}[t]{0.496\linewidth}
  \centering
  \includegraphics[width=\linewidth,height=5.7cm,keepaspectratio]{\WeightsFig}
\end{minipage}
\hfill
\begin{minipage}[t]{0.496\linewidth}
  \centering
  \includegraphics[width=\linewidth,height=5.7cm,keepaspectratio]{\EquityCurveFig}
\end{minipage}

\vspace{1pt}
\noindent
\begin{minipage}[t]{0.496\linewidth}
  \centering{\fontsize{7}{9}\selectfont\textcolor{CaptionGray}{GMV allocation across six assets on Oct 1, 2025.}}
\end{minipage}
\hfill
\begin{minipage}[t]{0.496\linewidth}
  \centering{\fontsize{7}{9}\selectfont\textcolor{CaptionGray}{Equity curve (Oct 1, 2025 -- Dec 31, 2025) with controlled drawdowns.}}
\end{minipage}

{\fontsize{8}{10.3}\selectfont\justifying
\textbf{Key insights.}
\begin{enumerate}
  \item INFY received the largest allocation (28.43\%) due to minimum variance contribution.
  \item The GMV solution remains diversified across all 6 stocks, with M\&M capped to 3.13\% due to volatility impact.
  \item The profile is conservative by design: variance minimization takes priority over return maximization.
  \item Sharpe ratio improved to 0.2113 (vs. 0.1628 baseline).
  \item Total return reached 7.33\% over the 3-month forward period (approximately 29.3\% annualized).
\end{enumerate}
}

% =========================
% Section 6
% =========================
\SectionBar{6 \;|\; Project Structure and Implementation}
{\fontsize{8}{10.3}\selectfont\justifying
\textbf{Directory structure.}
\begin{verbatim}
Assignment2_Project/
|- src/
|  |- data/
|  |  |- fetch_market_data.py      # Yahoo Finance OHLCV data
|  |  |- fetch_fundamentals.py     # Screener.in scraping + synthetic fallback
|  |  |- fetch_macro_data.py       # USD-INR, Oil, NIFTY, VIX, CPI
|  |  |- fetch_sentiment.py        # NewsAPI + FinBERT sentiment
|  |- features/
|  |  |- build_features.py         # Feature engineering pipeline
|  |- models/
|  |  |- train_model.py            # LSTM training with walk-forward CV
|  |- portfolio/
|  |  |- portfolio.py              # GMV optimization and backtesting
|- notebooks/
|  |- main.ipynb
|- data/processed/
|  |- market_daily.csv
|  |- fundamentals_daily.csv
|  |- fundamentals_quarterly.csv
|  |- macro_daily.csv
|  |- news_sentiment.csv
|  |- features.pkl
|  |- models.pkl
|  |- predictions.pkl
|  |- model_metrics.pkl
|  |- portfolio_results.pkl
|- reports/
|  |- report.md
|- requirements.txt
\end{verbatim}

\textbf{Pipeline execution flow (main.ipynb).}
\begin{enumerate}
  \item Step 1a: Fetch market data (8,916 rows).
  \item Step 1b: Fetch fundamentals data (13,152 rows).
  \item Step 1c: Fetch macro data (1,564 rows).
  \item Step 1d: Fetch sentiment data (13,152 rows).
  \item Step 2: Build features (58 features per ticker) and split train/test.
  \item Step 3: Train RNN models (6 models with early stopping).
  \item Step 4: Construct portfolio (GMV weights and performance metrics).
\end{enumerate}

\textbf{Technology stack.}
\begin{itemize}
  \item Data collection: \texttt{yfinance}, \texttt{requests}, \texttt{beautifulsoup4}, \texttt{newsapi-python}
  \item ML framework: \texttt{torch} (PyTorch)
  \item Portfolio optimization: \texttt{scipy.optimize} (SLSQP)
  \item Data processing: \texttt{pandas}, \texttt{numpy}, \texttt{scikit-learn}
  \item NLP/sentiment: \texttt{transformers} (FinBERT), \texttt{sentencepiece}
  \item Visualization: \texttt{matplotlib}, \texttt{seaborn}
  \item Utilities: \texttt{joblib}, \texttt{tqdm}, \texttt{jupyter}
\end{itemize}

\textbf{Data and modeling details.}
\begin{itemize}
  \item Market data: 1,486 trading days per ticker (Jan 2020 -- Dec 2025).
  \item Feature dimension: 58 features per ticker (lagged returns, technical, fundamentals, macro, sentiment).
  \item Train/test split: 1,376 training days and 60 test days.
  \item Sequence length: 30 days.
  \item LSTM: 2 layers, 64 hidden units, FC(32), dropout 0.3, weight decay $1\times10^{-5}$, gradient clipping 1.0.
  \item Validation: 5-fold expanding-window walk-forward.
  \item Optimization objective: minimum variance with long-only constraints, full investment, and max 50\% per asset.
  \item Covariance estimation: Ledoit-Wolf shrinkage.
  \item Rebalancing setup: single-period optimization on Oct 1, 2025.
\end{itemize}
}

% =========================
% Conclusion
% =========================
\SectionBar{Conclusion}
{\fontsize{8}{10.3}\selectfont\justifying
This project delivered an end-to-end multi-dimensional forecasting and portfolio management system for 6 major Indian equities.

\textbf{Project accomplishments.}
\begin{enumerate}
  \item Real fundamentals integration through screener.in scraping with fallback handling.
  \item Full ML pipeline from data acquisition to portfolio construction.
  \item Fusion of market, fundamental, macro, and sentiment data streams.
  \item Robust per-ticker LSTM training with walk-forward time-series validation.
  \item GMV optimization with shrinkage covariance and practical constraints.
\end{enumerate}

\textbf{Final performance summary.}
\begin{itemize}
  \item Sharpe ratio: 0.2113.
  \item Total return: 7.33\% over the 3-month forward window.
  \item Maximum drawdown: -2.52\%.
  \item Best single-stock model: HDFCBANK (MSE = 0.000049).
  \item Total daily observations: 13,152 across 6 assets.
\end{itemize}
}

\vspace{6pt}
\noindent\textcolor{MidGray}{\rule{\linewidth}{0.8pt}}
\vspace{2pt}
\begin{center}
  {\fontsize{6.5}{8}\selectfont\textcolor{FooterGray}{\FooterText}}
\end{center}

\end{document}
